\section{The QU-fitting technique}
Visualizing a model is one thing; connecting it to real data is another. QU-fitting does that by directly fitting a model's predicted fractional Stokes $q(\lambda)=Q(\lambda)/I(\lambda)$ and $u(\lambda)=U(\lambda)/I(\lambda)$ (see the Complex polarization help page) to the observed $q(\lambda)$, $u(\lambda)$, rather than fitting $p(\lambda)$ and $\chi(\lambda)$ directly. Measurement noise enters Stokes $I$, $Q$, $U$ as (to good approximation) independent Gaussians, and $q$/$u$ inherit that directly since they are just linear rescalings of $Q$/$U$ by $I$; $p=\sqrt{q^2+u^2}$ and $\chi=\frac12\arctan(u/q)$ do not, since both are non-linear combinations of $q$ and $u$, so fitting on $q$, $u$ keeps the statistics simple and avoids the bias a naive fit directly to $p$/$\chi$ would introduce.

Pressing "Load data" reads a csv file containing a frequency column and Stokes $I$, $Q$, $U$ (error columns are optional); polvista auto-detects the units from the column headers and normalizes every Stokes parameter by dividing by $I$ at the lowest available frequency. For $n$ loaded frequencies this gives $n$ complex data points $q_j+iu_j$, i.e. $2n$ real numbers to fit; the number of free parameters in the selected model is denoted $k$.

Fractional Stokes $q$, $u$ data alone cannot constrain a model's spectral-shape parameters ($\alpha$ and, for a two-component model, $\varepsilon$) since both cancel out of the fractional polarization by construction. Polvista therefore fixes $\varepsilon=0.5$ and estimates $\alpha$ directly from the loaded $I(\nu)$ data before either fitting method below ever runs: a closed-form power-law regression anchored through $I(\nu_{\rm min})=1$, or, under the SSA/Thermal spectral shapes, a nonlinear regression (optionally solving jointly for the turnover frequency $\nu_0$ too). This keeps the model's Stokes $I$ curve consistent with the real intensity spectrum without adding potentially correlated free parameters to the QU-fitting stage itself.

To ease convergence to physically meaningful values, every fit (least-squares or Bayesian) bounds the model's parameters: $p_0\in[0,70]\%$, $\phi\in[-5,5]\times10^6$ rad m$^{-2}$, $\sigma_\phi\in[0,5]\times10^6$ rad m$^{-2}$, $f\in[0,1]$, and $s\in[0,100]$. Because $\chi_0$ is a period-$\pi$ quantity, a direct bound on the EVPA can trap an optimizer against its $\pm\pi/2$ edges and produce an unreliable exploration of parameter space. Polvista avoids this by fitting every $(p_0,\chi_0)$ pair in an unconstrained cartesian parametrization instead, $a=p_0\cos(2\chi_0)$, $b=p_0\sin(2\chi_0)$, converting back to physical values only once a result is in hand: going all the way around the EVPA circle is then just walking around the origin in $(a,b)$ space, with no seam for the optimizer to get stuck against.

Following Anderson et al. (2016), the log-likelihood every fit in polvista aims to maximize is:
\begin{equation}\label{eq:loglike}
    \displaystyle \mathcal{L} = -\frac{1}{2}\sum_{j=1}^{n}\left[\left(\frac{q_j-\tilde q_j}{\sigma_{q_j}}\right)^2 + \left(\frac{u_j-\tilde u_j}{\sigma_{u_j}}\right)^2 + 2\ln\left(2\pi\sigma_{q_j}\sigma_{u_j}\right)\right]
\end{equation}
where $(\tilde q_j,\tilde u_j)$ are the model's predictions at frequency bin $j$. Both fitting paths below share this exact likelihood (including its normalization term), so the resulting $\chi^2$, $\ln\mathcal{L}$, AIC, and BIC values are directly comparable between them.
\\

\subsection{Simple least-squares}
Pressing "Fit!" under the Visualization tab minimizes $-\mathcal{L}$ with a robust (soft-L1) nonlinear least-squares solver, using the current slider values as the initial guess. Soft-L1 down-weights the influence of any single badly-discrepant point relative to an ordinary sum of squares, so one noisy channel is less likely to drag the whole fit off course. At the end of the fit, the sliders take the best-fit values and every plot updates to reflect that solution; the fit's $\chi^2$, reduced $\chi^2$, $\ln\mathcal{L}$, AIC, and BIC are reported alongside it.

Least-squares finds one local optimum from one starting point. If the true solution is genuinely ambiguous, with several distinct parameter combinations that explain the data comparably well, it will only ever report whichever one the initial guess happened to fall closest to, with no indication that the others exist.
\\

\subsection{Bayesian nested sampling}
The Sampling tab addresses that ambiguity directly, by mapping out the full posterior instead of chasing a single optimum. If pyMultiNest, a python interface to the Fortran MultiNest package (Feroz et al. 2009), is installed, this tab becomes usable once data is loaded. For a model with parameters $\theta$ and data $D$, the posterior probability distribution of the parameters follows from Bayes' theorem:
\begin{equation}\label{eq:bayes}
    \displaystyle \mathcal{P}(\theta|D) = \frac{\pi(\theta)\,\mathcal{P}(D|\theta)}{Z(D)}
\end{equation}
where $\pi(\theta)$ is the prior, $\log\mathcal{P}(D|\theta)=\mathcal{L}(D|\theta)$ is the log-likelihood of Equation \ref{eq:loglike}, and $Z(D)$ is the marginal likelihood, or Bayesian evidence:
\begin{equation}\label{eq:evidence}
    \displaystyle Z(D) = \int e^{\mathcal{L}(D|\theta)}\,\pi(\theta)\,d\theta
\end{equation}
obtained by integrating the likelihood over the entire prior rather than evaluating it at one point. Nested sampling (Skilling 2004) estimates $Z$ by replacing that high-dimensional integral with a one-dimensional one over the enclosed prior volume: starting from $n_{\rm live}$ points drawn from the full prior, each iteration discards the point with the lowest likelihood and replaces it with a new one drawn from the prior, restricted to the region of higher likelihood than the point just discarded. Each discarded point contributes a small weight to a running sum that converges to $Z$, while the surviving live points progressively converge onto the region(s) of high likelihood, i.e., the posterior itself.

The Sampling tab lets you set the output directory for the run, and MultiNest's own hyperparameters: the number of live points $n_{\rm live}$ (more live points resolve the posterior, and any multiple modes, more finely, at the cost of more likelihood evaluations), the sampling efficiency (MultiNest's target acceptance efficiency for its ellipsoidal sampling; lower values around $0.2-0.3$, the default, are recommended for a reliable evidence estimate, while higher values up to $\sim0.8$ trade that reliability for speed and suit quick parameter estimates better), and the evidence tolerance (the stopping criterion: sampling halts once the estimated remaining evidence contribution falls below it). Below these, one prior-bounds row is shown per parameter kind actually present in the model (skipping $\alpha$/$\varepsilon$, held fixed exactly as in the least-squares path, and EVPA, always sampled over its full range using the same seam-avoiding reparametrization idea as the cartesian trick above); $\phi$ and $\sigma_\phi$ share a "Log-uniform" toggle, recommended (and on by default) since their physically plausible range spans several decades and a plain linear prior over that range wastes almost all its mass at the largest magnitudes.

Because a given model can have several distinct solutions that fit the data comparably well, MultiNest can identify and characterize multiple candidate solutions at once, whenever its live points naturally diverge into disjoint, well-separated clusters during sampling. Once a run finishes, polvista merges any sufficiently nearby raw clusters into physically distinct families before reporting a result: two raw clusters are grouped into the same family if their marginal EVPA, $\phi$, and $\sigma_\phi$ estimates lie within a fixed significance threshold of one another (accounting for $\chi_0$'s $\pi$-periodicity); for models with an exact component-swap symmetry, this also checks the swapped labeling, so that true degeneracy collapses into one family instead of splitting its evidence between two apparently distinct modes. Each family's member clusters are then pooled, by evidence-weighted resampling, into one posterior sample set.

If distinct families still remain after that merge, polvista reports the one whose marginal-median point best explains the data (highest point log-likelihood), with the 16th and 84th percentiles of each parameter setting its asymmetric lower/upper uncertainty. This result is shown in a new "Corner plot" tab, a triangle plot of the model's parameters, alongside any alternative families that survived clustering with a non-negligible share of the evidence; every raw sample remains accessible through the files MultiNest itself saved to the chosen output directory, and "Load samples" can reload and re-cluster a previous run's output without sampling again.

Because small or sparsely-sampled datasets are especially prone to this kind of ambiguity, always check the Corner plot tab before trusting a single reported solution: an ideal one should carry high posterior mass (the fraction of samples clustered around it), a reduced $\chi^2$ near $1$, and look relatively simple in the resulting $p$/EVPA and Stokes $I$, $Q$, $U$ curves, with no more oscillation than the data actually demands.
